% ----- CHAUPAI 15 -----

\newpage

\subsection*{\deva{पञ्चदश चौपाई} (Fifteenth Chaupai)}
\addcontentsline{toc}{subsection}{\deva{पञ्चदश चौपाई} (Fifteenth Chaupai) — \deva{जमु कुबेर...}}

\begin{minipage}[t]{0.55\textwidth}
\vspace{0pt}

{\large \linespread{1.5}\selectfont
\deva{जमु कुबेर दिगपाल जहाँ ते।}\\
\deva{कबि कोबिद कहि सके कहाँ ते॥}
\par}

\vspace{0.5em}

{\large \linespread{1.5}\selectfont
\telu{జము కుబేర దిగపాల జహాఁ తే।}\\
\telu{కబి కోబిద కహి సకే కహాఁ తే॥}
\par}

\vspace{0.5em}

{\large \linespread{1.3}\selectfont
\textit{jamu kubera digapāla jahāṃ te |}\\
\textit{kabi kobida kahi sake kahāṃ te ||}
\par}
\end{minipage}%
\hfill
\begin{minipage}[t]{0.42\textwidth}
\vspace{0pt}
\begin{tcolorbox}[anchorbox]
\textbf{Remaining Grounded:} Even when surrounded by celestial beings, absolute authorities, and the masters of wealth and death who cannot measure his greatness, Hanuman remains entirely humble. True confidence doesn't require constant external validation, titles, or showing off.
\vspace{0.5em}\newline He carried immense, cosmic capabilities without ever letting his ego demand external validation from poets or scholars.\newline\null\hfill\href{https://www.valmiki.iitk.ac.in/sloka?field_kanda_tid=4&language=dv&field_sarga_value=66}{\textit{Kishkindha Kanda, 66:26-29}}
\end{tcolorbox}
\end{minipage}

\vspace{0.5em}
\subsubsection*{\textbf{Visual Rhythm Map:}}
\begin{table}[h!]
\centering
\renewcommand{\arraystretch}{1.2}
\setlength{\tabcolsep}{0pt}
\begin{tabularx}{\textwidth}{|*{16}{>{\centering\arraybackslash\hsize=1.0\hsize}X|}}
\hline
\multicolumn{2}{|>{\centering\arraybackslash\hsize=2.0\hsize}X|}{\textbf{\laghu{ज}\laghu{मु}} \par (2)} &
\multicolumn{4}{>{\centering\arraybackslash\hsize=4.0\hsize}X|}{\textbf{\laghu{कु}\guru{बे}\laghu{र}} \par (4)} &
\multicolumn{5}{>{\centering\arraybackslash\hsize=5.0\hsize}X|}{\textbf{\laghu{दि}\laghu{ग}\guru{पा}\laghu{ल}} \par (5)} &
\multicolumn{3}{>{\centering\arraybackslash\hsize=3.0\hsize}X|}{\textbf{\laghu{ज}\guru{हाँ}} \par (3)} &
\multicolumn{2}{>{\centering\arraybackslash\hsize=2.0\hsize}X|}{\textbf{\guru{ते}} \par (2)} \\
\hline
\end{tabularx}
\vspace{-1.5em}
\end{table}

\begin{table}[h!]
\centering
\renewcommand{\arraystretch}{1.2}
\setlength{\tabcolsep}{0pt}
\begin{tabularx}{\textwidth}{|*{16}{>{\centering\arraybackslash\hsize=1.0\hsize}X|}}
\hline
\multicolumn{2}{|>{\centering\arraybackslash\hsize=2.0\hsize}X|}{\textbf{\laghu{क}\laghu{बि}} \par (2)} &
\multicolumn{4}{>{\centering\arraybackslash\hsize=4.0\hsize}X|}{\textbf{\guru{को}\laghu{बि}\laghu{द}} \par (4)} &
\multicolumn{2}{>{\centering\arraybackslash\hsize=2.0\hsize}X|}{\textbf{\laghu{क}\laghu{हि}} \par (2)} &
\multicolumn{3}{>{\centering\arraybackslash\hsize=3.0\hsize}X|}{\textbf{\laghu{स}\guru{के}} \par (3)} &
\multicolumn{3}{>{\centering\arraybackslash\hsize=3.0\hsize}X|}{\textbf{\laghu{क}\guru{हाँ}} \par (3)} &
\multicolumn{2}{>{\centering\arraybackslash\hsize=2.0\hsize}X|}{\textbf{\guru{ते}} \par (2)} \\
\hline
\end{tabularx}
\vspace{-1.5em}
\end{table}

\vspace{1em}
\textbf{ \deva{सरल-संस्कृतम्}:}\\
 \deva{यत्र यमराजः, कुबेरः, दिक्पालाः च}\\
 \deva{(तव यशः वर्णयितुं न शक्नुवन्ति), ततः कथं कवयः कोविदाः (पण्डिताः) च तत् वक्तुं शक्नुवन्ति?}\\


When Yama (God of Death), Kubera (God of Wealth), and the Guardians of the Directions\\
fail to describe your glory, then how can mere poets and scholars ever describe it?


\vspace{1em}
\newpage
\textbf{\deva{प्रतिपदार्थः} (Meaning of each word):}
\begin{table}[h!]
\centering
\renewcommand{\arraystretch}{1.2}
\begin{tabularx}{\textwidth}{@{} l l X X @{}}
\toprule
\textbf{Awadhi} & \textbf{Sanskrit} & \textbf{English} & \textbf{Grammar \& Notes} \\
\midrule
\deva{जमु} / \telu{జము}  &  \deva{यमः}  &  God of Death  & \\ 
\deva{कुबेर} / \telu{కుబేర}  &  \deva{कुबेरः}  &  God of Wealth  &   \\
\deva{दिगपाल} / \telu{దిగపాల}  &  \deva{दिक्पालाः}  &  Direction Guardians  & $k \rightarrow g$ shift \\ 
\deva{जहाँ ते} / \telu{జహాఁ తే}  &  \deva{यत्र / यस्मात्}  &  Where from / Even they  &   \\
\deva{कबि} / \telu{కబి}  &  \deva{कविः}  &  Poet  & \\ 
\deva{कोबिद} / \telu{కోబిద}  &  \deva{कोविदः (पण्डितः)}  &  Scholar  & \\ 
\deva{कहि सके} / \telu{కహి సకే}  &  \deva{वक्तुं शक्नुवन्ति}  &  Can say/describe  &   \\
\deva{कहाँ ते} / \telu{కహాఁ తే}  &  \deva{कुतः / कथम्}  &  How / From where  &   \\
\bottomrule
\end{tabularx}
\end{table}

\newpage
